\section{导论}
\begin{frame}{工程层次}
\begin{center}
\resizebox{!}{6.6cm}{\begin{tikzpicture}[font=\small,>=Latex,
  rung/.style={draw,thick,rounded corners,align=center,inner sep=4pt,minimum height=.8cm,font=\footnotesize},
  platform/.style={draw=darkorange!85!black,line width=1.3pt,rounded corners,align=center,inner sep=5pt,fill=darkorange!14,minimum width=11.4cm,minimum height=1.5cm},
  pillar/.style={draw,thick,rounded corners,align=center,inner sep=4pt,minimum width=5cm,minimum height=1.35cm}]
\node[rung,fill=black!3,draw=black!35,text=black!60,minimum width=4.6cm] (p) at (0,-.35) {\textbf{提示词}工程};
\node[rung,fill=black!6,draw=black!45,text=black!70,minimum width=5.6cm] (c) at (0,.95) {\textbf{上下文}工程};
\draw[->,very thick,black!45] (p)--(c);
\draw[dashed,black!30] (3,-.92)--(3,1.58);
\node[anchor=west,font=\footnotesize,text=black!55,align=left,text width=3cm] at (3.15,.33) {专题 2：\\提示词与上下文};
\node[platform] (h) at (0,2.95) {\textbf{\normalsize Harness 工程}\\[2pt]\footnotesize Hook $\cdot$ 检查关卡 $\cdot$ 契约 $\cdot$ 溯源 $\cdot$ 状态 $\cdot$ 脚本};
\draw[->,very thick,black!70] (c)--(h);
\node[pillar,fill=steelblue!12,draw=steelblue!80!black,anchor=south] (l) at (-2.9,4.1) {\textbf{循环（Loop）工程}\\\footnotesize 时间 $\cdot$ 状态 $\cdot$ 停止规则};
\node[pillar,fill=forestgreen!10,draw=forestgreen!70!black,anchor=south] (g) at (2.9,4.1) {\textbf{图（Graph）工程}\\\footnotesize 拓扑结构 $\cdot$ 谁能看到什么};
\draw[line width=2pt,darkorange!85!black] (-2.9,3.7)--(-2.9,4.1);
\draw[line width=2pt,darkorange!85!black] (2.9,3.7)--(2.9,4.1);
\draw[darkred,line width=1.1pt,rounded corners=5pt] (-6.05,1.98) rectangle (6.05,5.80);
\node[fill=white,inner sep=4pt,text=darkred,font=\small\bfseries] at (0,5.80) {本讲：Harness、Loop 与 Graph};
\node[font=\footnotesize,text=black!65,align=center] at (0,6.4) {Loop 与 Graph 建立在 Harness 之上，而非更高一层的阶梯。};
\end{tikzpicture}
}
\end{center}
\end{frame}
\begin{frame}{构建智能体量化系统}
\begin{center}\resizebox{.94\textwidth}{!}{\begin{tikzpicture}[
  box/.style={draw=blue,fill=blue!5,rounded corners=2pt,text width=3.3cm,minimum height=1.5cm,align=center,font=\small,inner sep=5pt},
  link/.style={-{Stealth[length=5pt]},thick,draw=black!60}]
\node[box] (n0) at (1.7,0) {\textbf{交易想法}\\带日期的数据};
\node[box] (n1) at (6.6,0) {\textbf{策略 + 代码}\\经核验的结果};
\draw[link] (n0) -- (n1);
\node[box,draw=forestgreen,fill=forestgreen!6] (n2) at (11.5,0) {\textbf{独立审查}\\接受、修订或停止};
\draw[link] (n1) -- (n2);
\draw[link,draw=darkred,rounded corners=3pt] (n2.south) -- (11.5,-1.55) -- node[below,font=\footnotesize,text=darkred]{修订} (6.6,-1.55) -- (n1.south);
\end{tikzpicture}
}\end{center}
\begin{itemize}
\item 把交易想法与带日期的数据转化为有证据支撑的策略。
\item 返回证据与局限；验收需要可追溯、可复算、可恢复的工作。
\end{itemize}
\end{frame}
\begin{frame}[t]{完整架构}
\centering
\resizebox{!}{6.0cm}{\begin{tikzpicture}[
  lead/.style={draw,rounded corners=2pt,line width=.7pt,align=center,font=\footnotesize,inner sep=4pt},
  worker/.style={draw,dashed,rounded corners=2pt,fill=white,text width=1.30cm,minimum height=.62cm,align=center,font=\scriptsize,inner sep=2pt},
  flow/.style={-{Stealth[length=4pt]},line width=.65pt,draw=black!60},
  edge/.style={font=\scriptsize,align=center,text=black!65,fill=white,inner sep=1pt}]
\draw[darkorange!85!black,line width=1.1pt,rounded corners=5pt]
  (-.45,.45) rectangle (15.75,-5.65);
\node[fill=white,inner sep=4pt,font=\small\bfseries,text=darkorange!85!black]
  at (7.65,.45) {演进：经评估的系统改动};
\node[lead,draw=forestgreen,fill=forestgreen!6,text width=4.5cm,minimum height=.68cm] (o) at (7.65,-.20)
  {\textbf{Orchestrator}\\分派 + 状态 + 路由};
\node[font=\footnotesize,align=center,text width=3.1cm] (idea) at (1.5,.85)
  {\textbf{用户想法}\\任务授权 + 限额};
\node[font=\scriptsize,align=center,text=darkred,text width=3.4cm] (human) at (13.6,.85)
  {\textbf{人工制定规则}\\并保留批准权};
\draw[flow] (idea.south)--(1.5,-.20)--(o.west);
\draw[flow,draw=darkred] (human.south)--(13.6,-.20)--(o.east);

\draw[blue,line width=.7pt,rounded corners=3pt,fill=blue!4] (0,-1.0) rectangle (5.4,-3.7);
\draw[darkred,line width=.7pt,rounded corners=3pt,fill=red!4] (7.8,-1.0) rectangle (15.3,-3.7);
\node[font=\footnotesize\bfseries,text=blue] at (2.7,-1.28) {Generator 负责人};
\node[font=\footnotesize\bfseries,text=darkred] at (11.55,-1.28) {Critic 负责人};
\node[font=\scriptsize,text=black!65] at (2.7,-1.60) {协调构建专项 Agent};
\node[font=\scriptsize,text=black!65] at (11.55,-1.60) {每轮开展全新审查};
\node[worker,draw=blue] (gd) at (1.0,-2.30) {数据};
\node[worker,draw=blue] (gf) at (2.7,-2.30) {特征};
\node[worker,draw=blue] (gb) at (4.4,-2.30) {回测};
\node[worker,draw=darkred,text width=1.52cm] (cd) at (8.8,-2.30) {数据与证据};
\node[worker,draw=darkred,text width=1.52cm] (cr) at (10.65,-2.30) {稳健性};
\node[worker,draw=darkred,text width=1.52cm] (ce) at (12.5,-2.30) {执行};
\node[worker,draw=darkred,text width=1.52cm] (ck) at (14.35,-2.30) {风险与控制};
\node[lead,draw=blue,fill=white,text width=4.3cm,minimum height=.52cm] (gi) at (2.7,-3.18)
  {\textbf{整合 + 测试}};
\node[lead,draw=darkred,fill=white,text width=6.4cm,minimum height=.52cm] (cm) at (11.55,-3.18)
  {\textbf{合并 + 覆盖检查}};
\foreach \w in {gd,gf,gb}{\draw[flow] (2.7,-1.83)--(\w.north);\draw[flow] (\w.south)--(gi.north);}
\foreach \w in {cd,cr,ce,ck}{\draw[flow] (11.55,-1.83)--(\w.north);\draw[flow] (\w.south)--(cm.north);}
\draw[flow] (o.south west)--(4.7,-.82)--(2.7,-.82)--(2.7,-1.0);
\draw[flow] (o.south east)--(10.3,-.82)--(11.55,-.82)--(11.55,-1.0);
\node[edge,anchor=south] at (2.7,-.75) {分派};
\node[edge,anchor=south] at (11.55,-.75) {检查通过后};
\draw[flow,draw=blue] (5.4,-2.0)--(7.8,-2.0);
\node[edge,text=blue] at (6.6,-1.62) {可执行检查关卡\\经由 Orchestrator};
\node[lead,draw=forestgreen,fill=forestgreen!6,text width=4.6cm,minimum height=.65cm] (packet) at (7.65,-4.30)
  {\textbf{Orchestrator：修复任务包}\\问题项 + 决策：继续 / 结束};
\draw[flow] (cm.south)--(11.55,-4.30)--(packet.east);
\node[edge,anchor=west] at (11.7,-4.1) {问题项交回\\Orchestrator};
\draw[flow] (packet.west)--(2.7,-4.30)--(2.7,-3.7);
\node[edge,anchor=east] at (2.55,-4.1) {继续时\\发回修复};
\node[draw=black!35,rounded corners=2pt,fill=gray!5,text width=14.4cm,align=center,
  minimum height=.45cm,font=\scriptsize,inner sep=4pt] at (7.65,-5.15)
  {\textbf{共享状态与证据：}请求 ID、产物、哈希值、问题项与版本};
\end{tikzpicture}
}
\par\raggedright
\smallskip
\begin{itemize}
\item 演进在各次运行之间评估系统改动；最终批准权仍由人掌握。
\end{itemize}
\end{frame}
\begin{frame}{金融：主要差异}
\begin{itemize}
\item \textbf{重大后果}
\begin{itemize}
\item 错误可能影响客户、资金和机构义务，部分操作难以撤销。
\item 支持结论的证据本身不构成行动授权。
\end{itemize}
\item \textbf{时间信息}
\begin{itemize}
\item 金融结果取决于决策时点可得的信息与产生该结果的版本。
\item 应区分重建已发生的运行、复算固定计算和新的模型运行。
\end{itemize}
\item \textbf{边界授权}
\begin{itemize}
\item 需要明确谁能行动、授权范围、限额以及谁能停止工作流。
\item 研究验收、版本批准和资金使用是不同决定。
\end{itemize}
\item \textbf{持续运行与受控变更}
\begin{itemize}
\item 中断后恢复已核验工作；评估改动时保留既有控制与旧版本。
\end{itemize}
\end{itemize}
\end{frame}

\begin{frame}{职责分工}
\begin{center}
\resizebox{.96\textwidth}{!}{\begin{tikzpicture}[
  human/.style={draw=darkred,fill=red!4,rounded corners=3pt,line width=.9pt,
    text width=4.4cm,minimum height=1.25cm,align=center,font=\small,inner sep=5pt},
  task/.style={draw=steelblue,fill=blue!3,rounded corners=2pt,line width=.7pt,
    text width=1.95cm,minimum height=1.3cm,align=center,font=\footnotesize,inner sep=4pt},
  arrow/.style={-{Stealth[length=5pt]},line width=.8pt,draw=black!60},
  edge/.style={font=\scriptsize,text=black!65,fill=white,inner sep=2pt}]
\node[anchor=east,font=\scriptsize\bfseries,text=darkred] at (1.2,0) {人};
\node[human] (judgment) at (5.0,0) {\textbf{判断}\\初始想法 + 系统设计};
\node[human] (understanding) at (10.55,0) {\textbf{决策}\\资金使用};
\fill[gray!8,rounded corners=5pt] (1.6,-1.8) rectangle (14.0,-4.2);
\node[anchor=east,align=center,font=\scriptsize\bfseries,text=steelblue] at (1.2,-2.8) {量化\\系统};
\node[task] (task0) at (2.95,-2.55) {\textbf{数据}\\带日期的输入};
\node[task] (task1) at (5.37,-2.55) {\textbf{设计}\\信号 + 代码};
\node[task] (task2) at (7.79,-2.55) {\textbf{测试}\\收益 + 成本};
\node[task] (task3) at (10.21,-2.55) {\textbf{审查}\\质疑 + 修复};
\node[task] (task4) at (12.629999999999999,-2.55) {\textbf{报告}\\证据};
\node[draw=black!35,fill=white,rounded corners=2pt,text width=11.5cm,minimum height=.5cm,font=\footnotesize,align=center,inner sep=3pt] at (7.8,-3.65) {代码：检查、路由与状态保存};
\draw[arrow] (judgment.south) -- node[left,edge]{任务授权} (5.0,-1.8);
\draw[arrow] (10.55,-1.8) -- node[right,edge]{证据} (understanding.south);
\end{tikzpicture}
}
\end{center}
\medskip
\begin{itemize}
\item 人选择初始想法与系统设计，再决定如何使用资金。
\end{itemize}
\end{frame}
\begin{frame}{提纲}
\begin{itemize}
\item \textbf{长程工作流设计}
\begin{itemize}
\item Harness：运行环境与控制
\item Loop：循环与停止
\item Graph：依赖与交接
\item Evolution：经评估的系统改动
\end{itemize}
\item \textbf{量化系统案例}
\begin{itemize}
\item 构建 QT 组件与检查
\item 演进一个可复用的系统版本
\end{itemize}
\end{itemize}
\end{frame}
\begin{frame}{Orchestrator、Generator 与 Critic}
\begin{center}
\resizebox{.94\textwidth}{!}{\begin{tikzpicture}[
  part/.style={draw=black!30,fill=gray!5,rounded corners=2pt,text width=4.0cm,minimum height=1.0cm,align=center,font=\footnotesize,inner sep=4pt},
  card/.style={draw=blue,fill=blue!4,rounded corners=3pt,text width=5.4cm,minimum height=1.35cm,align=center,font=\small,inner sep=7pt}]
\node[part,text=black!45] at (1.9,0) {1. 导论};
\node[part,draw=darkred,fill=red!5,text=darkred] at (6.55,0) {\textbf{2. 长程工作流设计}};
\node[part,text=black!45] at (11.2,0) {3. 案例：量化系统};
\node[card] at (3,-1.7) {\textbf{Harness}\\运行环境与控制};
\node[card] at (10.1,-1.7) {\textbf{循环}\\重复执行与停止};
\node[card] at (3,-3.5) {\textbf{图}\\依赖关系与任务交接};
\node[card,draw=forestgreen,fill=forestgreen!5] at (10.1,-3.5) {\textbf{演进}\\经过评估的系统改动};
\end{tikzpicture}
}
\end{center}
\end{frame}
\begin{frame}{为何分开角色？}
\begin{center}
\resizebox{.94\textwidth}{!}{\begin{tikzpicture}[
  role/.style={draw,rounded corners=3pt,text width=3.8cm,minimum height=1.15cm,align=center,font=\small,inner sep=6pt},
  flow/.style={-{Stealth[length=5pt]},thick,draw=black!60},
  edge/.style={font=\footnotesize,fill=white,inner sep=2pt}]
\node[role,draw=forestgreen,fill=forestgreen!5] (o) at (6.6,0) {\textbf{Orchestrator}\\确定范围、委派、路由};
\node[role,draw=blue,fill=blue!4] (g) at (2,-2.0) {\textbf{Generator}\\提议、实现、测试};
\node[role,draw=darkred,fill=red!4] (c) at (11.2,-2.0) {\textbf{Critic}\\质疑、核验、报告};
\node[draw=black!50,fill=gray!6,rounded corners=2pt,text width=2.4cm,minimum height=.8cm,align=center,font=\footnotesize] (gate) at (6.6,-2.0) {\textbf{代码检查关卡}\\经由 Orchestrator};
\draw[flow] (o.west)-|node[above,edge,pos=.3]{请求}(g.north);
\draw[flow] (g)--(gate);
\draw[flow] (gate)--(c);
\draw[flow] (c.north)|-node[above,edge,pos=.7]{问题项}(o.east);
\node[draw=black!35,fill=gray!5,rounded corners=2pt,text width=12.5cm,align=center,font=\footnotesize,inner sep=6pt] at (6.6,-3.6)
  {Orchestrator 发出修复请求，或记录终止结果。};
\end{tikzpicture}
}
\end{center}
\begin{itemize}
\item 主会话协调工作；控制器代码检查每一次状态转换。
\item Generator 和 Critic 都不能自行扩大权限或绕过检查。
\end{itemize}
\end{frame}
\begin{frame}{角色契约}
\begin{center}
\small
\begin{tabular}{@{}P{4.2cm}P{8.3cm}@{}}
\toprule
\textbf{设计问题} & \textbf{职责分离} \\
\midrule
产出方自行验收 & Critic 在单独的上下文中检查证据 \\
\rowcolor{gray!10}
审查改变了目标 & 任务授权与验收标准保持不变 \\
控制权归属不清 & Orchestrator 决定路由；代码核验状态转换 \\
\bottomrule
\end{tabular}

\end{center}
\begin{itemize}
\item 独立的上下文有助于 Critic 质疑方案，但两个 Agent 仍可能犯同样的错误。
\item 在任务委派和独立核验确有价值时采用这些角色。
\end{itemize}
\end{frame}
\begin{frame}{每个 Harness 内部}
\begin{center}
\resizebox{.92\textwidth}{!}{\begin{tikzpicture}[
  box/.style={draw=blue,fill=blue!5,rounded corners=2pt,text width=3.1cm,minimum height=.9cm,align=center,font=\footnotesize,inner sep=4pt},
  flow/.style={-{Stealth[length=4pt]},line width=.7pt,draw=black!60},
  label/.style={font=\scriptsize,align=center,fill=white,inner sep=1pt,text=black!60}]
\draw[forestgreen,line width=.8pt,rounded corners=4pt,fill=forestgreen!3] (0,.35) rectangle (13.2,-4.90);
\node[font=\small\bfseries,text=forestgreen] at (6.6,-.05) {Harness：受控的运行环境};
\node[box] (context) at (2.0,-1.2) {\textbf{(1) 上下文}\\指令 + 带日期的证据};
\node[box,fill=white] (agent) at (6.6,-1.2) {\textbf{Agent}\\规划下一步};
\node[box,draw=black!60,fill=gray!8] (tools) at (11.2,-1.2) {\textbf{(2) 工具}\\权限 + 计算};
\node[box,draw=forestgreen,fill=forestgreen!6] (state) at (2.0,-3.25) {\textbf{(3) 状态}\\ID、哈希值\\未解决问题项};
\node[box,draw=darkred,fill=red!4] (gate) at (6.6,-3.25) {\textbf{(4) 检查}\\结构约束 + 证据};
\node[box,draw=black!60,fill=gray!6] (result) at (11.2,-3.25) {\textbf{交回负责人}\\产物或错误};
\draw[flow] (context)--(agent);
\draw[flow] (agent)--node[above,label]{行动}(tools);
\draw[flow] (tools.south)--(11.2,-2.2)--node[below,label]{观察结果}(8.2,-2.2)--(agent.south east);
\draw[flow] (agent.south)--node[left,label]{提交}(gate.north);
\draw[flow] (gate)--node[above,label]{通过}(result);
\draw[flow,draw=darkred] (gate.south)--(6.6,-4.10)--node[below,label,text=darkred]{失败：返回错误}(11.2,-4.10)--(result.south);
\draw[flow,dashed] (gate.west)--node[above,label]{记录}(state.east);
\draw[flow,dashed] (state.north)--node[left,label]{只读\\上下文}(context.south);
\node[font=\scriptsize,align=center,text=black!60] at (6.6,-4.62) {工具调用、重试次数与总用时均有上限};
\end{tikzpicture}
}
\end{center}
\begin{itemize}
\item 四个编号组件共同构成每个角色的 Harness；控制器拥有运行状态并负责推进流程。
\end{itemize}
\end{frame}
\begin{frame}{契约与边界}
\begin{itemize}
\item \textbf{任务与产物契约}
\begin{itemize}
\item 明确任务、固定输入、权限、返回格式及完成/失败条件。
\item 每份产物都关联请求及其支持证据。
\end{itemize}
\item \textbf{权限边界}
\begin{itemize}
\item 披露文件和网页是数据，不是权限来源。
\item 付费数据请求仍需核验用途、供应商、金额、有效期、身份和批准人。
\item 课堂 QT 没有支付或券商权限。
\end{itemize}
\end{itemize}
\end{frame}
\begin{frame}{Hook 与状态控制器}
\begin{center}\resizebox{.82\textwidth}{!}{\begin{tikzpicture}[
  box/.style={draw=blue,fill=blue!5,rounded corners=2pt,text width=3.3cm,minimum height=1.5cm,align=center,font=\small,inner sep=5pt},
  link/.style={-{Stealth[length=5pt]},thick,draw=black!60}]
\node[box] (n0) at (1.7,0) {\textbf{生命周期事件}\\工具调用结束};
\node[box] (n1) at (6.6,0) {\textbf{Hook 脚本}\\自动记录事件};
\draw[link] (n0) -- (n1);
\node[box,draw=forestgreen,fill=forestgreen!6] (n2) at (11.5,0) {\textbf{控制器门控}\\核验案例证据};
\draw[link] (n1) -- (n2);
\end{tikzpicture}
}\end{center}
\begin{itemize}
\item \textbf{Hook}
\begin{itemize}
\item 宿主事件发生时运行回调，记录事件元数据。
\item 只有宿主支持事件阻断时，才可能阻止操作。
\end{itemize}
\item \textbf{状态控制器}
\begin{itemize}
\item 核验提交产物、哈希、问题项和数值检查。
\item 决定运行是前进、重试、等待还是停止。
\end{itemize}
\item Hook 观察事件；控制器管理工作流状态。
\end{itemize}
\end{frame}
\begin{frame}{状态与检查点}
\begin{center}\resizebox{.82\textwidth}{!}{\begin{tikzpicture}[box/.style={draw=blue,fill=blue!4,rounded corners=2pt,text width=5.5cm,minimum height=1.25cm,align=center,font=\small,inner sep=6pt}]
\node[box] at (3,0) {已完成\\输出已经核验};
\node[box] at (10,0) {待完成\\请求已经发出};
\node[box,draw=forestgreen,fill=forestgreen!5] at (3,-2) {状态不明\\无法确认是否完成};
\node[box,draw=forestgreen,fill=forestgreen!5] at (10,-2) {下一步\\先检查，再继续};
\end{tikzpicture}
}\end{center}
\begin{itemize}
\item (1) 身份：案例、请求、候选方案和提示词/版本 ID。
\item (2) 进度：待处理角色、轮次、重试次数和截止时间。
\item (3) 证据：用于发现产物变化的路径与哈希。
\item (4) 问题项：未解决问题和已核验修复。
\item 每次状态转换时保存已核验状态；对话最后一行不能说明哪些步骤已经完成。
\end{itemize}
\end{frame}
\begin{frame}{审计证据}
\begin{center}\resizebox{.94\textwidth}{!}{\begin{tikzpicture}[
  box/.style={draw=blue,fill=blue!5,rounded corners=2pt,text width=2.6cm,minimum height=1.5cm,align=center,font=\small,inner sep=5pt},
  link/.style={-{Stealth[length=5pt]},thick,draw=black!60}]
\node[box] (n0) at (1.5,0) {\textbf{数据来源}\\带日期快照 ID};
\node[box] (n1) at (4.9,0) {\textbf{数值计算}\\落盘产物文件};
\draw[link] (n0) -- (n1);
\node[box] (n2) at (8.3,0) {\textbf{独立审查}\\缺陷台账与路由};
\draw[link] (n1) -- (n2);
\node[box,draw=forestgreen,fill=forestgreen!6] (n3) at (11.7,0) {\textbf{版本决策}\\版本号与审批记录};
\draw[link] (n2) -- (n3);
\end{tikzpicture}
}\end{center}
\begin{itemize}
\item 将 JSON 证据保存在本地，并保留 Git 版本历史。每次批准都记录批准人、时间戳，以及获准的操作或版本。
\item 另一位分析师可以据此重建可观察的操作与文件；记录并不揭示模型隐藏的推理过程。
\end{itemize}
\end{frame}
